\documentclass[11pt,letterpaper]{article}
\usepackage[margin=1in]{geometry}
\usepackage{graphicx} % Required for inserting images
\usepackage{tikz}
\usepackage{amsthm,amsmath,amssymb}
\usepackage{dsfont}
\usepackage{algorithm}
\usepackage{algpseudocode}
\usepackage{xcolor}
\usepackage{soul}
\usepackage{enumitem}

\usepackage[hidelinks]{hyperref}
\usepackage{cleveref}
\usepackage{etoolbox}

% Theorems
\newtheorem{theorem}{Theorem}[section]
\newtheorem{corollary}[theorem]{Corollary}
\newtheorem{lemma}[theorem]{Lemma}
\newtheorem{claim}[theorem]{Claim}
\theoremstyle{definition}
\newtheorem{definition}[theorem]{Definition}
\newtheorem{remark}[theorem]{Remark}
\newtheorem{observation}[theorem]{Observation}

\crefname{theorem}{Theorem}{Theorems}
\Crefname{theorem}{Theorem}{Theorems}
\crefname{corollary}{Corollary}{Corollaries}
\Crefname{corollary}{Corollary}{Corollaries}
\crefname{lemma}{Lemma}{Lemmas}
\Crefname{lemma}{Lemma}{Lemmas}
\crefname{claim}{Claim}{Claims}
\Crefname{claim}{Claim}{Claims}
\crefname{definition}{Definition}{Definitions}
\Crefname{definition}{Definition}{Definitions}
\crefname{remark}{Remark}{Remarks}
\Crefname{remark}{Remark}{Remarks}
\crefname{observation}{Observation}{Observations}
\Crefname{observation}{Observation}{Observations}

% The environments above share the `theorem` counter, and cleveref names a label
% after its counter -- so without these aliases every \cref to a lemma, claim,
% etc. comes out as "Theorem". \crefalias is local to the environment's group,
% so each one is typed correctly while the shared numbering is left alone.
\AtBeginEnvironment{corollary}{\crefalias{theorem}{corollary}}
\AtBeginEnvironment{lemma}{\crefalias{theorem}{lemma}}
\AtBeginEnvironment{claim}{\crefalias{theorem}{claim}}
\AtBeginEnvironment{definition}{\crefalias{theorem}{definition}}
\AtBeginEnvironment{remark}{\crefalias{theorem}{remark}}
\AtBeginEnvironment{observation}{\crefalias{theorem}{observation}}

\newcommand{\reals}{\mathbb{R}}
\newcommand{\Span}{\mathrm{span}}
\newcommand{\coreq}{\mathrm{coreq}}
\newcommand{\cone}{\mathrm{cone}}

\newcommand{\argmax}{\operatorname*{arg\,max}}
\newcommand{\argmin}{\operatorname*{arg\,min}}
\newcommand{\tr}{\operatorname{tr}}

\newcommand{\calA}{\mathcal{A}}
\newcommand{\calB}{\mathcal{B}}
\newcommand{\calC}{\mathcal{C}}
\newcommand{\calD}{\mathcal{D}}
\newcommand{\calE}{\mathcal{E}}
\newcommand{\calF}{\mathcal{F}}
\newcommand{\calG}{\mathcal{G}}
\newcommand{\calH}{\mathcal{H}}
\newcommand{\calI}{\mathcal{I}}
\newcommand{\calJ}{\mathcal{J}}
\newcommand{\calK}{\mathcal{K}}
\newcommand{\calL}{\mathcal{L}}
\newcommand{\calM}{\mathcal{M}}
\newcommand{\calN}{\mathcal{N}}
\newcommand{\calO}{\mathcal{O}}
\newcommand{\calP}{\mathcal{P}}
\newcommand{\calQ}{\mathcal{Q}}
\newcommand{\calR}{\mathcal{R}}
\newcommand{\calS}{\mathcal{S}}
\newcommand{\calT}{\mathcal{T}}
\newcommand{\calU}{\mathcal{U}}
\newcommand{\calV}{\mathcal{V}}
\newcommand{\calW}{\mathcal{W}}
\newcommand{\calX}{\mathcal{X}}
\newcommand{\calY}{\mathcal{Y}}
\newcommand{\calZ}{\mathcal{Z}}

\newcommand{\bfa}{\mathbf{a}}
\newcommand{\bfb}{\mathbf{b}}
\newcommand{\bfc}{\mathbf{c}}
\newcommand{\bfd}{\mathbf{d}}
\newcommand{\bfe}{\mathbf{e}}
\newcommand{\bff}{\mathbf{f}}
\newcommand{\bfg}{\mathbf{g}}
\newcommand{\bfh}{\mathbf{h}}
\newcommand{\bfi}{\mathbf{i}}
\newcommand{\bfj}{\mathbf{j}}
\newcommand{\bfk}{\mathbf{k}}
\newcommand{\bfl}{\mathbf{l}}
\newcommand{\bfm}{\mathbf{m}}
\newcommand{\bfn}{\mathbf{n}}
\newcommand{\bfo}{\mathbf{o}}
\newcommand{\bfp}{\mathbf{p}}
\newcommand{\bfq}{\mathbf{q}}
\newcommand{\bfr}{\mathbf{r}}
\newcommand{\bfs}{\mathbf{s}}
\newcommand{\bft}{\mathbf{t}}
\newcommand{\bfu}{\mathbf{u}}
\newcommand{\bfv}{\mathbf{v}}
\newcommand{\bfw}{\mathbf{w}}
\newcommand{\bfx}{\mathbf{x}}
\newcommand{\bfy}{\mathbf{y}}
\newcommand{\bfz}{\mathbf{z}}

\newcommand{\cost}{\mathrm{COST}}
\newcommand{\ALG}{\mathrm{ALG}}
\newcommand{\OPT}{\mathrm{OPT}}
\newcommand{\Flow}{\textsc{Flow}}
\newcommand{\ouralg}{\textsc{LineTracker}}

\newcommand{\EE}{\mathbb{E}}
\newcommand{\Var}{\mathrm{Var}}
\newcommand{\Conv}{\mathrm{Conv}}
\newcommand{\One}{\mathds{1}}

\newcommand{\interval}[1]{\lfloor #1 \rceil}
\newcommand{\dis}[1]{#1}
\newcommand{\off}{\mathrm{off}}
\newcommand{\on}{\mathrm{on}}
\newcommand{\org}{\mathrm{org}}
\newcommand{\rel}{\mathrm{rel}}
\newcommand{\new}{\mathrm{new}}

\newcommand{\solution}{\medskip\noindent\textbf{Solution.}\par\medskip}

\title{47-851: Modern Convex Optimization \\ Homework 3}
\author{Poon Tze-Yang}
\date{November 2026}


\begin{document}

\maketitle

\noindent Lecturer: Karan Singh

\begin{enumerate}[series=q]
\item \textbf{Grads and Dists (10 points)}
\begin{enumerate}[label=\textbf{Part \Alph*.}, leftmargin=*]
\item Consider a function $f : \reals^d \to \reals$ that is $\mu$-strongly convex and $L$-smooth. Let $x^* \in \argmin_x f(x)$. Prove for all $x \in \reals^d$ that
\[
\frac{1}{2L}\|\nabla f(x)\|^2 \leq f(x) - f(x^*) \leq \frac{1}{2\mu}\|\nabla f(x)\|^2.
\]
\item Prove that $f$ is $\mu$-strongly convex if and only if for all $x, y \in \reals^d$, we have
\[
(\nabla f(x) - \nabla f(y))^\top (x - y) \geq \mu\|x - y\|^2.
\]
\end{enumerate}
\end{enumerate}

\solution
A.
By definition, $f$ is $\mu$-strongly convex iff
\[
f(y) \geq f(x) + \nabla f(x)^\top (y-x) + \frac{\mu}{2}\|y-x\|^2.
\]
Rearranging:
\[
f(y) -f(x) \geq \nabla f(x)^\top (y-x) + \frac{\mu}{2}\|y-x\|^2.
\]
Taking the minimum of both sides w.r.t. $y$:
\[
\min_y\{f(y) -f(x)\} \geq \min_y\{\nabla f(x)^\top (y-x) + \frac{\mu}{2}\|y-x\|^2\}.
\]
The left hand side is equal to $f(x^*)-f(x)$.
The right hand side is minimized when the derivative with respect to $y$ is equal to 0:
\begin{align*}
    \nabla f(x) + \frac{\mu}{2}(2y-2x) = 0\quad\Rightarrow\quad y -x = -\frac{1}{\mu}\nabla f(x).
\end{align*}
Substituting back in:
\[
\min_y\{\nabla f(x)^\top (y-x) + \frac{\mu}{2}\|y-x\|^2\} = -\frac{1}{2\mu}\|\nabla f(x)\| ^2.
\]
Combining both sides:
\[
f(x^*) - f(x) \geq -\frac{1}{2\mu}\|\nabla f(x)\| ^2 \quad\Rightarrow\quad f(x) - f(x^*)  \leq \frac{1}{2\mu}\|\nabla f(x)\| ^2.
\]

A function $f$ is $L$-smooth if
\[
f(y)\leq f(x) + \nabla f(x)^\top(y-x) + \frac{L}{2}\left\|(y-x)\right\|^2
\]
Essentially the same argument as above gives 
\[
f(x)-f(x^*) \geq \frac{1}{2L}\|\nabla f(x)\|^2.
\]
\medskip
B.
For the forward direction, use $\mu$-strong convexity at $x,y$:
\begin{align*}
    f(y) - f(x) \geq \nabla f(x)^\top (y-x) + \frac{\mu}{2}\|y-x\|^2. \\
    f(x) - f(y) \geq \nabla f(y)^\top (x-y) + \frac{\mu}{2}\|y-x\|^2.
\end{align*}
Combining the two inequalities:
\begin{align*}
    \nabla f(y)^\top (y-x) - \frac{\mu}{2}\|y-x\|^2 \geq \nabla f(x)^\top (y-x) + \frac{\mu}{2}\|y-x\|^2 \\
    \Rightarrow (\nabla f(x) - \nabla f(y))^\top (x - y) \geq \mu\|x - y\|^2.
\end{align*}

For the reverse direction:

We integrate the directional derivative from $y$ to $x$:
\begin{align*}
    f(x) - f(y) &= \int_{0}^{\|x-y\|}\nabla f(y+\frac{t}{\|x-y\|}(x-y))^\top(\frac{1}{\|x-y\|}(x-y))dt \\
    &= \int_{0}^{\|x-y\|}\frac{1}{t}\left(\nabla f(y+\frac{t}{\|x-y\|}(x-y))^\top(\frac{t}{\|x-y\|}(x-y)) \right)dt
\end{align*}
Rearranging the hypothesis with $x\gets y+(t/\|x-y\|) (x-y), y\gets y$, we get:
\begin{align*}
    \nabla f(y+\frac{t}{\|x-y\|}(x-y))^\top(\frac{t}{\|x-y\|}(x-y)) \geq \nabla f(y)^\top(\frac{t}{\|x-y\|}(x-y)) + \mu\left\|\frac{t}{\|x-y\|}(x-y)\right\|^2.
\end{align*}
Hence, 
\begin{align*}
    f(x) - f(y) &\geq \int_{0}^{\|x-y\|} \nabla f(y)^\top(\frac{t}{\|x-y\|}(x-y)) + \mu\left\|\frac{t}{\|x-y\|}(x-y)\right\|^2 dt \\
    & = \int_{0}^{\|x-y\|} \frac{1}{t}\left(\nabla f(y)^\top(\frac{t}{\|x-y\|}(x-y)) + \mu\left\|\frac{t}{\|x-y\|}(x-y)\right\|^2 \right)dt \\
    & = \int_{0}^{\|x-y\|} \nabla f(y)^\top(\frac{1}{\|x-y\|}(x-y)) + \mu t \ dt \\
    & = \nabla f(y)^\top(x-y) + \frac{\mu}{2} \|x-y\|^2 .\\
\end{align*}
This is the definition for strong convexity.


\begin{enumerate}[resume=q]
\item \textbf{Just Enough Convexity (10 points)}

Consider a $L$-smooth function $f : \reals^d \to \reals$ that satisfies for all $x$ that
\[
f(x) - \min_{x^*} f(x^*) \leq \frac{1}{2\mu}\|\nabla f(x)\|^2.
\]
\begin{enumerate}[label=\textbf{Part \Alph*.}, leftmargin=*]
\item Prove that such a function may not even be convex (much less strongly convex).
\item Nevertheless, prove that gradient descent on such a function, with an appropriate step size, produces iterates so that
\[
f(x_t) - \min_{x^*} f(x^*) \leq O\bigl(e^{-\mu \frac{t}{L}}\bigr).
\]
\end{enumerate}
\end{enumerate}

\solution

For a univariate function $f$, from Taylor's theorem, $f(y) = f(x) + f'(x)(x-y) + \frac{1}{2}f''(\xi)(x-y)^2$ for some $\xi\in [x,y]$.
If we can bound the absolute value of $f''(\xi)$ for $\xi\in \reals$, we have shown that $f$ is $L$-smooth for appropriate $L$.

Consider the function $f(x) = x^2 + 3\sin^2x$. We have that $f'(x) = 2x + 3\sin 2x$ and $f''(x) = 2 + 6 \cos 2x $.
Since $f''(x)\in [-4,8]$, $f$ is $L$-smooth for appropriate $L$.



\begin{enumerate}[resume=q]
\item \textbf{Small Grad Bounds (10 points)}
\begin{enumerate}[label=\textbf{Part \Alph*.}, leftmargin=*]
\item Fix a convex set $K \subset \reals^d$ of diameter $D$. Recall that online gradient guarantees the following regret bound against a sequence of adversarially chosen convex $G$-Lipschitz functions $f_t$.
\[
\mathrm{Regret} = \sum_{t=1}^T f_t(x_t) - \min_{x^* \in K} \sum_{t=1}^T f_t(x^*) \leq O\bigl(GD\sqrt{T}\bigr),
\]
With suitable dynamically chosen step sizes, show that one can guarantee that
\[
\mathrm{Regret} \leq O\left(D\sqrt{\sum_{t=1}^T \|\nabla f_t(x_t)\|^2}\right).
\]

\emph{Hint:} Set $\eta_t \propto \frac{1}{\sqrt{\sum_{s=1}^t \|\nabla f_s(x_s)\|^2}}$.
\item Assume that all loss functions $f_t : \reals^d \to \reals_{\geq 0}$ are in addition non-negative, well-defined and $L$-smooth over all of $\reals^d$. Our online algorithm must still play points in $K$ and the regret is still defined with respect to the best point in $K$. Prove that the same algorithm guarantees
\[
\mathrm{Regret} \leq O\left(\sqrt{LD^2\left(\min_{x^* \in K} \sum_{t=1}^T f_t(x^*)\right)} + LD^2\right).
\]
This means whenever there is a point with zero net loss, we get constant regret, independent of $T$.

\emph{Hint:} Use part A from the first question.
\end{enumerate}
\end{enumerate}

\solution



\end{document}
